\section*{अध्याय \ref{ch:b3:petrological-thermodynamics} --- चट्टानों की ऊष्मागतिकी}

\begin{solution}{exo:b3:petrological-thermodynamics:1}
$\Delta_rH = -3029.3 - 910.7 + 3935.0 = \qty{-5.0}{kJ}$; $\Delta_rS = 133.5 +
41.5 - 207.4 = \qty{-32.4}{J.K^{-1}}$; $\Delta_rV = 60.40 + 22.69 - 100.07 =
\qty{-16.98}{cm^3} = \qty{-1.698}{J.bar^{-1}}$। ढलान $32.4/1.698 =
\qty{19.1}{bar/K}$। \qty{873}{K} पर: $P = 1 + (873.15\times(-32.4) +
5000)/(-1.698) = \qty{13.7}{kbar}$।
\end{solution}

\begin{solution}{exo:b3:petrological-thermodynamics:2}
$c = 2$ और $p = 3 > c$: नियम भंग होता है। तीनों प्रावस्थाएँ केवल
\ce{Mg2SiO4 + SiO2 <=> 2 MgSiO3} की एकचरी रेखा पर ही सह-अस्तित्व रखती हैं;
यादृच्छिक रूप से चुनी गई कोई चट्टान वहाँ नहीं होती, इसलिए यह समुच्चय एक अपूर्ण अभिक्रिया का अभिलेख है,
उदाहरण के लिए, एंस्टैटाइट के आवरण से कवचबद्ध फ़ॉर्स्टेराइट, जिसने उसे क्वार्ट्ज़ से
अलग रखा।
\end{solution}

\begin{solution}{exo:b3:petrological-thermodynamics:3}
$a_{\text{ग्रॉसुलर}} = 0.20^3 = 0.008$; $a_{\text{पाइरोप}} = 0.15^3 =
\num{3.4e-3}$; $a_{\text{एनॉर्थाइट}} = 0.40$।
\end{solution}

\begin{solution}{exo:b3:petrological-thermodynamics:4}
$T_c = 15\,000/(2\times8.314) = \qty{902}{K}$; $W = 2\times8.314\times1000
= \qty{16.6}{kJ.mol^{-1}}$।
\end{solution}

\begin{solution}{exo:b3:petrological-thermodynamics:5}
$\Delta_rG = -5432.2 + 2\times(-2088.8) - 730.2 + 3\times3478.1 =
\qty{94.3}{kJ}$; $P = 1 + 94\,300/6.61 = \qty{14.3}{kbar}$। प्रायोगिक
रेखा $22.80\times726.85 - 1093 = \qty{15.5}{kbar}$ देती है,
\qty{1.2}{kbar} अधिक।
\end{solution}

\begin{solution}{exo:b3:petrological-thermodynamics:6}
$RT = \qty{7483}{J.mol^{-1}}$। $x = 0.5$ पर: $7483\ln0.5 + 20\,000/4 = -5187 +
5000 = \qty{-187}{J.mol^{-1}}$। $x = 0.111$ पर: $7483\times(0.111\ln0.111 +
0.889\ln0.889) + 20\,000\times0.111\times0.889 = -2609 + 1974 =
\qty{-635}{J.mol^{-1}}$, और समरूपता से 0.889 पर भी यही मान। दोनों पार्श्वों की समान मात्राओं का
स्थूल संघटन 0.5 है, और गिब्स ऊर्जा
\qty{-635}{J.mol^{-1}} है: विभाजन से यह \qty{448}{J.mol^{-1}} घट जाती है।
\end{solution}

\begin{solution}{exo:b3:petrological-thermodynamics:7}
$c + 2 = 5$ प्रावस्थाएँ और 5 रेखाएँ। (D) $A + C = 2B$: $0 + 2/3 = 2\times1/3$। (A)
$B + D = 2C$: $1/3 + 1 = 2\times2/3$। (C) $2A + D = 3B$: $0 + 1 = 3\times1/3$।
(B) $A + 2D = 3C$: $0 + 2 = 3\times2/3$।
\end{solution}

\begin{solution}{exo:b3:petrological-thermodynamics:8}
लीवर नियम से, 0.889 पर की प्रावस्था $(0.3 - 0.111)/(0.889 - 0.111) =
0.243$, \qty{24}{\percent} बनाती है, और 0.111 पर की प्रावस्था \qty{76}{\percent}।
\end{solution}

\begin{solution}{exo:b3:petrological-thermodynamics:9}
$\partial P/\partial X_{\text{ऐन}} = -3RT/(6.21X_{\text{ऐन}}) =
-3\times8.314\times900/(6.21\times0.30) = \qty{-12000}{bar}$ प्रति इकाई
$X_{\text{ऐन}}$: 0.01 की त्रुटि दाब को \qty{120}{bar} खिसका देती है।
\end{solution}

\begin{solution}{exo:b3:petrological-thermodynamics:10}
$G'_{\text{मिश्र}} = RT\ln[x/(1 - x)] + W(1 - 2x)$ और $G''_{\text{मिश्र}} =
RT/[x(1 - x)] - 2W$, जहाँ $x(1 - x) = RT/(2W)$ शून्य है। $x = 0.3$ के लिए: $T =
2\times20\,000\times0.21/8.314 = \qty{1010}{K}$; \omterm{def:b3:petrological-thermodynamics:solvus}{सॉल्वस}
$20\,000\times0.4/(8.314\ln(7/3)) = \qty{1136}{K}$ पर पहुँचता है। दोनों के बीच क्रिस्टल
उपस्थायी होता है: विभाजन इसकी ऊर्जा घटा देता, परंतु एक छोटा
उतार-चढ़ाव इसे बढ़ा देता है, इसलिए दूसरी प्रावस्था को केंद्रक बनाना होगा।
\end{solution}

\begin{solution}{exo:b3:petrological-thermodynamics:11}
चूँकि $G_{\text{मिश्र}}(x) = G_{\text{मिश्र}}(1 - x)$, उभयनिष्ठ स्पर्श-रेखा
वक्र को दो सममित बिंदुओं $x$ और $1 - x$ पर स्पर्श करती है, और क्षैतिज है:
$G'_{\text{मिश्र}}(x) = 0$, जो यही समीकरण है। हल करने पर, $T(x) = W(1 -
2x)/[R\ln((1 - x)/x)]$। $x = 1/2 - \varepsilon$ लेने पर, $1 - 2x = 2\varepsilon$
और $\ln((1/2 + \varepsilon)/(1/2 - \varepsilon)) \approx 4\varepsilon$, अतः
$T \to W/(2R)$ जब $\varepsilon \to 0$; और $x < 1/2$ के लिए $T(x)$ बढ़ता है, जैसे-जैसे $x$
1/2 की ओर पहुँचता है; इसलिए यह सीमा अधिकतम है।
\end{solution}

\begin{solution}{exo:b3:petrological-thermodynamics:12}
$(2.0 - 0.08\times10)/0.92 = \qty{1.30}{wt\percent}$। मैंगनीज़ गार्नेट के पक्ष में है:
\omterm{def:b3:petrological-thermodynamics:pseudosection}{प्रभावी स्थूल संघटन} में इसकी कम मात्रा होने पर परिधियाँ भिन्न क्षेत्र में
और भिन्न आइसोप्लेथ के अनुदिश बढ़ती हैं, जैसा संपूर्ण-चट्टान विश्लेषण से
अपेक्षित होता है।
\end{solution}

\begin{solution}{pb:b3:petrological-thermodynamics:1}
\textbf{1.} $\Delta_rH^\circ = -6640.0 + 2\times(-2593.8) - 910.7 + 3\times4234.0
= \qty{-36.3}{kJ}$; $\Delta_rS^\circ = 260.1 + 165.6 + 41.5 - 597.9 =
\qty{-130.7}{J.K^{-1}}$; $\Delta_rV = 125.28 + 88.30 + 22.69 - 302.37 =
\qty{-66.10}{cm^3} = \qty{-6.61}{J.bar^{-1}}$।

\textbf{2.} $-36\,300 + 298.15\times130.7 = \qty{+2.67}{kJ}$: धनात्मक है, अतः
एनॉर्थाइट सतह पर स्थायी है।

\textbf{3.} $P = 1 + (900\times(-130.7) + 36\,300)/(-6.61) = \qty{12.3}{kbar}$;
ढलान $130.7/6.61 = \qty{19.8}{bar/K}$।

\textbf{4.} $1 + 81\,500/6.61 = \qty{12.3}{kbar}$ \qty{900}{K} पर; $1 +
119\,800/6.61 = \qty{18.1}{kbar}$, \qty{1200}{K} पर।

\textbf{5.} 626.85 और \qty{926.85}{\celsius} पर प्रायोगिक रेखा
13.2 और \qty{20.0}{kbar} देती है: सारणी-मानों से 0.9 और \qty{1.9}{kbar} ऊपर।

\textbf{6.} $1914\times6.61 = \qty{12.7}{kJ}$। संयुक्त अनिश्चितता
$\sqrt{3.2^2 + (3\times2.0)^2 + (2\times2.0)^2 + 1.0^2} = \qty{8.0}{kJ}$ है: अंतर इसका
1.6 गुना है, जो ऊष्मामिति के लिए संभव है, और यही कारण है कि यह रेखा
प्रयोगों से ली गई है।

\textbf{7.} $\Delta_rC_p$ ठोसों में छोटा है, और इसके दोनों समाकल आंशिक रूप से
एक-दूसरे को निरस्त करते हैं: सारणी-मान रेखा को केवल 26 और \qty{112}{bar} से खिसकाते हैं।

\textbf{8.} \ce{SiO2}, \ce{Al2O3}, \ce{FeO}, \ce{MgO}, \ce{MnO}, \ce{CaO},
\ce{Na2O}, \ce{K2O} और \ce{H2O}: $c = 9$। प्रावस्थाएँ: गार्नेट, बायोटाइट,
मस्कोवाइट, प्लैजियोक्लेज़, काइनाइट, क्वार्ट्ज़ और द्रव, $p = 7$।

\textbf{9.} $F = 9 - 7 + 2 = 4$: दाब, तापमान और दो संघटनात्मक
चर। $p = 7 \le c = 9$: नियम लागू होता है, और यह संयोजन
दाब--तापमान तल के किसी क्षेत्र पर स्थायी रह सकता है।

\textbf{10.} साम्य में स्थित दो शुद्ध \ce{Al2SiO5} बहुरूप
काइनाइट = सिलिमेनाइट एकचरी रेखा पर स्थित होते हैं: $P$ और $T$ के बीच
एक संबंध देती है, और तब एक तापमापी से दाब ज्ञात हो जाता है।

\textbf{11.} द्रव जल का रासायनिक विभव आरोपित करता है: 6 ठोस प्रावस्थाओं
के लिए 8 मुक्त घटक हैं, $F = 8 - 6 + 2 = 4$, और $p \le c$ अब भी
लागू होता है।

\textbf{12.} $a_{\text{ग्रॉसुलर}} = 0.10^3 = 0.001$, $a_{\text{एनॉर्थाइट}} =
0.30$; $\ln K = \ln(0.001/0.30^3) = -3.30$.

\textbf{13.} $K_D = (0.12/0.74)/(0.45/0.55) = 0.1622/0.8182 = 0.198$.

\textbf{14.} $T = (52\,635 + 2488)/(20.20 + 24.94\times1.618) = 55\,123/60.57
= \qty{910}{K}$, या \qty{637}{\celsius}।

\textbf{15.} $P_0 = 22.80\times636.9 - 1093 = \qty{13429}{bar}$; $(3\times8.314
\times910.1/6.21)\ln(1/3) = \qty{-4016}{bar}$; $P = \qty{9.41}{kbar}$.

\textbf{16.} \qty{9.41}{kbar} पर \qty{644.2}{\celsius}, फिर
\qty{9.55}{kbar}; उस दाब पर \qty{644.9}{\celsius} और
\qty{9.56}{kbar}, इसके बाद अपरिवर्तित: \qty{645}{\celsius} और
\qty{9.56}{kbar}.

\textbf{17.} \qty{641}{\celsius} और \qty{8.79}{kbar}। सारणी-मान
अंत्य-सदस्य रेखा को लगभग \qty{1}{kbar} कम दर्शाते हैं (भाग I);
प्रायोगिक रेखा रखी गई है।

\textbf{18.} $X_{\text{ऐन}} = 0.35$ के लिए: \qty{642}{\celsius} और
\qty{8.93}{kbar}, \qty{0.63}{kbar} कम।

\textbf{19.} \qty{918}{K} पर: $P = 1 + (918\times12.6 - 7700)/0.571 =
\qty{6.8}{kbar}$, \qty{9.56}{kbar} से कम: काइनाइट, जैसा प्रेक्षित हुआ।

\textbf{20.} $z = \num{9.56e8}/(2800\times9.81) = \qty{34.8}{km}$.

\textbf{21.} $\Delta P = (3\times8.314\times918/6.21)\ln(0.07/0.10) =
\qty{-1.3}{kbar}$: क्रोड लगभग \qty{4.8}{km} उथले स्तर पर विकसित हुआ।
गार्नेट दफ़न के दौरान बढ़ा, और इसका सीमांत सबसे गहरे बिंदु को अभिलिखित करता है।

\textbf{22.} वार्डन शिस्ट ने चरम दाब
\textbf{\qty[text-series-to-math=true]{9.6}{kbar}} का
\qty{645}{\celsius} पर प्राप्त किया, जो लगभग \qty{35}{km} के दफ़न से मेल खाता है।
\end{solution}
